% Options for packages loaded elsewhere
\PassOptionsToPackage{unicode}{hyperref}
\PassOptionsToPackage{hyphens}{url}
\documentclass[
]{article}
\usepackage{xcolor}
\usepackage[margin=1in]{geometry}
\usepackage{amsmath,amssymb}
\setcounter{secnumdepth}{-\maxdimen} % remove section numbering
\usepackage{iftex}
\ifPDFTeX
  \usepackage[T1]{fontenc}
  \usepackage[utf8]{inputenc}
  \usepackage{textcomp} % provide euro and other symbols
\else % if luatex or xetex
  \usepackage{unicode-math} % this also loads fontspec
  \defaultfontfeatures{Scale=MatchLowercase}
  \defaultfontfeatures[\rmfamily]{Ligatures=TeX,Scale=1}
\fi
\usepackage{lmodern}
\ifPDFTeX\else
  % xetex/luatex font selection
\fi
% Use upquote if available, for straight quotes in verbatim environments
\IfFileExists{upquote.sty}{\usepackage{upquote}}{}
\IfFileExists{microtype.sty}{% use microtype if available
  \usepackage[]{microtype}
  \UseMicrotypeSet[protrusion]{basicmath} % disable protrusion for tt fonts
}{}
\makeatletter
\@ifundefined{KOMAClassName}{% if non-KOMA class
  \IfFileExists{parskip.sty}{%
    \usepackage{parskip}
  }{% else
    \setlength{\parindent}{0pt}
    \setlength{\parskip}{6pt plus 2pt minus 1pt}}
}{% if KOMA class
  \KOMAoptions{parskip=half}}
\makeatother
\usepackage{longtable,booktabs,array}
\usepackage{calc} % for calculating minipage widths
% Correct order of tables after \paragraph or \subparagraph
\usepackage{etoolbox}
\makeatletter
\patchcmd\longtable{\par}{\if@noskipsec\mbox{}\fi\par}{}{}
\makeatother
% Allow footnotes in longtable head/foot
\IfFileExists{footnotehyper.sty}{\usepackage{footnotehyper}}{\usepackage{footnote}}
\makesavenoteenv{longtable}
\usepackage{graphicx}
\makeatletter
\newsavebox\pandoc@box
\newcommand*\pandocbounded[1]{% scales image to fit in text height/width
  \sbox\pandoc@box{#1}%
  \Gscale@div\@tempa{\textheight}{\dimexpr\ht\pandoc@box+\dp\pandoc@box\relax}%
  \Gscale@div\@tempb{\linewidth}{\wd\pandoc@box}%
  \ifdim\@tempb\p@<\@tempa\p@\let\@tempa\@tempb\fi% select the smaller of both
  \ifdim\@tempa\p@<\p@\scalebox{\@tempa}{\usebox\pandoc@box}%
  \else\usebox{\pandoc@box}%
  \fi%
}
% Set default figure placement to htbp
\def\fps@figure{htbp}
\makeatother
\setlength{\emergencystretch}{3em} % prevent overfull lines
\providecommand{\tightlist}{%
  \setlength{\itemsep}{0pt}\setlength{\parskip}{0pt}}
\usepackage{float}
\usepackage{booktabs}
\usepackage{longtable}
\usepackage{array}
\usepackage{multirow}
\usepackage{wrapfig}
\usepackage{colortbl}
\usepackage{pdflscape}
\usepackage{tabu}
\usepackage{threeparttable}
\usepackage{threeparttablex}
\usepackage[normalem]{ulem}
\usepackage{makecell}
\usepackage{xcolor}
\usepackage{bookmark}
\IfFileExists{xurl.sty}{\usepackage{xurl}}{} % add URL line breaks if available
\urlstyle{same}
\hypersetup{
  pdftitle={Stats 272 Project: Final Report},
  pdfauthor={Aria Whalen and Gwynnie Hayes},
  hidelinks,
  pdfcreator={LaTeX via pandoc}}

\title{Stats 272 Project: Final Report}
\author{Aria Whalen and Gwynnie Hayes}
\date{}

\begin{document}
\maketitle

\subsection{Introduction}\label{introduction}

In order to identify unique species, researchers will often use a
variety of methods, including size, color, behavior, and sound. Size
measurements can be especially important when a specimen is visually
difficult to differentiate from other species or deceased.
Identification of species is crucial for conservation, as different
species and families will require a unique habitat for their lifestyle
and simply put, we cannot even begin to try and conserve a species if we
don't know what they are. Frog populations have been declining for
decades (Wake, 1991), often even without a specific known cause.
Utilizing body measurements of previously deceased frog specimens we
seek to build a model that can identify a species based off a handful of
measurements. Body size measurements are also known to be associated
with niche partioning in a variety of other animals, and a body size
based model could help explore which body parts are notably different
between families.

We chose to explore 2 frogs families colloquially known to contain
vastly different physical characteristics and habitats: Hylidae (tree
frogs) and Bufonidae (true toads). Frogs in the family of Hylidae are
known to live up in the trees of forested environments, though some live
in ponds and swamps. They are known to be long and skinny, and have
distinct toe pads that assist with climbing and strong legs for jumping
among the trees (AmphibiaWeb). Frogs in the family of Bufonidae are
typically known to be terrestrial/ground dwelling, they dig burrows and
prefer moist shady environments. (Bufonidae).

Both families have been studied extensively as focuses of work in
conservation, and are known to use their habitats very specifically.
Sanchez et. al (2010) found that the micro-habitats of \emph{Rhinella
fernandezae} (a species in the Bufonidae family) are not randomly
chosen, but a highly specific combination of soil hardness and broad
leaf plant cover. The choice of soil hardness is especially linked to
their tendancy for burrowing. Paris (2000) found that the number of
\emph{Litoria pearsoniana} (a species in the Hylidae family) were highly
correlated with the size of the local stream and presence of palm trees,
the latter of which is a limiting factor which is actively threatened by
forest clear cutting for cattle grazing. These papers show that even
small changes to an environment can have massive effects on the number
of species in an area.

\begin{figure}
\centering
\includegraphics[width=3.64583in,height=\textheight,keepaspectratio]{/home/rstudio/users/whalen4/salmon stomachs/Stats 2 Project/data/fig-1-full.png}
\caption{A: Snout-Vent Length (SVL), B: Lower Arm Length, C: Lower Leg
Length (Tibia), Eye Length, Head Length --- Guayasamin et al.,
2022}\label{id}
\end{figure}

\begin{longtable}[]{@{}
  >{\raggedright\arraybackslash}p{(\linewidth - 2\tabcolsep) * \real{0.1610}}
  >{\raggedright\arraybackslash}p{(\linewidth - 2\tabcolsep) * \real{0.8390}}@{}}
\caption{Description of Relevant Variables}\tabularnewline
\toprule\noalign{}
\begin{minipage}[b]{\linewidth}\raggedright
Measurement
\end{minipage} & \begin{minipage}[b]{\linewidth}\raggedright
Description
\end{minipage} \\
\midrule\noalign{}
\endfirsthead
\toprule\noalign{}
\begin{minipage}[b]{\linewidth}\raggedright
Measurement
\end{minipage} & \begin{minipage}[b]{\linewidth}\raggedright
Description
\end{minipage} \\
\midrule\noalign{}
\endhead
\bottomrule\noalign{}
\endlastfoot
Lower Leg Length & The length of the tibia and fibula, measured from the
heel to the proximal edge of the knee \\
Lower Arm Length & Measured from the wrist to the proximal end of the
elbow \\
Head Length & Measured from the jaw joint to the tip of the snout \\
Eye Length & Measured from the anterior corned to the posterior corner
of the eyelids \\
Naris-Snout Length & Measured from the Anterior edge of the naris to the
edge of the snout on the same side of the body \\
X Length Ratio & X Variable above was divided by Snout-Vent Length
(below) to remove size colinearity bias \\
Snout-Vent Length & Measured from the tip of the snout to vent \\
Family & Hylidae: Tree Frogs and their allies, Bufonidae: True Toads \\
\end{longtable}

\subsection{Materials and Methods}\label{materials-and-methods}

These data were collected by researchers at the University of Michigan
Museum of Zoology (UMMS) Division of Reptiles and Amphibians, from
specimens in their collection. The collection includes 2600 individual
specimens that contain 757 different species from 52 different families.
The frogs were measured in millimeters using digital calipers to
determine length and width of each body parts. Each row represents a
unique frog and only adult frogs were measured. Some species are
represented as many as 9 times, though many individuals have only 5 or
less observations, with this in mind we chose to represent the data on
the family level.

To explore these data for our research question (below) we had to do
some data manipulation to put it in a workable state. First, the data
were filtered to just include Hylidae and Bufonidae families and any
NA's were removed, which lowered the data to 898 full rows. Families
were also renamed to their common names, while Hylidae is technically a
family of Tree frogs \textbf{and their allies}, we just chose to name
them tree frogs for ease of understandability. Then, to remove size
based colinearity all variables were divided by SVL; as a result, all
explanatory variables in the project are described as ``relative to body
size'', and unit-less. For our logistic tests, we also multiplied the
explanatory variables by 20 to display them as a 5\% ratio to body size,
to account for small ratio values. Summary statitics are a median (to
account for high amount of outliers) with standard deviation to define
spread. Lastly, we used this data to identify an optimal generalized
logistic model using a stepwise analysis in both directions, which uses
a two-sided p-value.

\subsection{Results}\label{results}

\begin{longtable}[]{@{}lrr@{}}
\caption{Number of rows and \% of dataset by Family}\tabularnewline
\toprule\noalign{}
family & n & Percent \\
\midrule\noalign{}
\endfirsthead
\toprule\noalign{}
family & n & Percent \\
\midrule\noalign{}
\endhead
\bottomrule\noalign{}
\endlastfoot
Tree Frogs & 536 & 59.7 \\
True Toads & 362 & 40.3 \\
\end{longtable}

\begin{longtable}[t]{lrrrr}
\caption{\label{tab:unnamed-chunk-6}Summary of All Body Measurements for 
                                     Frogs + Toads}\\
\toprule
\multicolumn{1}{c}{ } & \multicolumn{2}{c}{Toads} & \multicolumn{2}{c}{Frogs} \\
\cmidrule(l{3pt}r{3pt}){2-3} \cmidrule(l{3pt}r{3pt}){4-5}
Measurements & Median & Std. Dev. & Median & Std. Dev.\\
\midrule
Tibio Fibula Length & 23.5 & 10.3 & 19.5 & 9.6\\
Antebrachial Length & 15.0 & 6.3 & 7.6 & 3.4\\
Head Length & 18.8 & 8.8 & 12.9 & 6.3\\
Eye Length & 5.4 & 2.5 & 3.8 & 1.6\\
Naris Snout Length & 1.7 & 1.0 & 1.4 & 0.7\\
\addlinespace
Snout-Vent Length & 59.3 & 27.5 & 37.5 & 17.5\\
Head Width & 20.2 & 11.1 & 12.0 & 6.2\\
Femur Length & 22.8 & 10.9 & 18.0 & 9.1\\
Inter-Oribital Length & 8.1 & 4.6 & 6.6 & 3.8\\
Inter-Narial Distance & 2.4 & 1.9 & 2.3 & 1.3\\
\addlinespace
Eye Naris Length Length & 3.9 & 1.9 & 3.4 & 2.2\\
\bottomrule
\end{longtable}

Table 3 defines the summary of all variables that were looked at in the
dataset. The first 6 rows are those identified as useful to the model by
the stepwise analysis, in order of importance. The remaining 5 variables
were marked as not statistically significant to the model. Though not a
rule, variables whose SD is less than half of the median are more
significant, and those where the SD is over half of the median are less
or in-significant.

\begin{longtable}[t]{lrrrr}
\caption{\label{tab:unnamed-chunk-7}Summary of All Body Measurements (but SVL) for Frogs + Toads \textbf{as a Ratio of Body Size}}\\
\toprule
\multicolumn{1}{c}{ } & \multicolumn{2}{c}{Toads} & \multicolumn{2}{c}{Frogs} \\
\cmidrule(l{3pt}r{3pt}){2-3} \cmidrule(l{3pt}r{3pt}){4-5}
Measurements & Median & Std. Dev. & Median & Std. Dev.\\
\midrule
Tibio Fibula Length Ratio & 0.391 & 0.053 & 0.526 & 0.056\\
Antebrachial Length Ratio & 0.248 & 0.031 & 0.198 & 0.024\\
Head Length Ratio & 0.314 & 0.034 & 0.342 & 0.034\\
Eye Length Ratio & 0.092 & 0.014 & 0.100 & 0.013\\
Naris Snout Length Ratio & 0.031 & 0.011 & 0.038 & 0.012\\
\addlinespace
Head Width Ratio & 0.342 & 0.043 & 0.315 & 0.039\\
Femur Length Ratio & 0.389 & 0.057 & 0.473 & 0.057\\
Inter-Oribital Length Ratio & 0.143 & 0.022 & 0.181 & 0.031\\
Inter-Narial Distance Ratio & 0.042 & 0.018 & 0.064 & 0.018\\
Eye Naris Length Length Ratio & 0.068 & 0.012 & 0.090 & 0.019\\
\bottomrule
\end{longtable}

Table 4 displays all variables divided by svl as a ratio of body size.
These numbers are decimals and unitless, and represent a similar trend
to the un-ratioed variables.

\begin{figure}[H]

{\centering \includegraphics{frog_project_files/figure-latex/unnamed-chunk-9-1} 

}

\caption{Boxplot of all Ratioed Variables, seperated by family. Most pairs have IQR overalap, though medians are statistically different in all plots (all p-values 1.4e-13 or lower). SVL represented in corner due to scale inequality, measured in millimeters -- all other measurements unitless.}\label{fig:unnamed-chunk-9}
\end{figure}

\begin{figure}[H]
\includegraphics{frog_project_files/figure-latex/unnamed-chunk-10-1} \caption{Density plots for all six variables in the model, showing individual probabilities of a specimin being a frog or toad across varying body measurement lengths}\label{fig:unnamed-chunk-10}
\end{figure}

There are noticable differences between the families on an individual
variable basis for all 6 measurements (Fig 1 + 2). The ratios of lower
leg length and lower arm length especially, show clear distinctions
between the two groups in both the boxplots and density curves. The
ratio of head length shows a similar density plot to Eye length, though
its more widely spread boxplot differentiates it. The ratios of eye
length and naris-snout length show density plots that are not nearly as
clean, and boxplots with medians that are very close together,
explaining why they are less useful to the model. Lastly, though SVL
displays a large difference in its scatterplot, the density plot shows
that the variable is not as clear cut as one might assume.

\begin{figure}
\centering
\pandocbounded{\includegraphics[keepaspectratio]{/home/rstudio/users/whalen4/salmon stomachs/Stats 2 Project/data/plotly.png}}
\caption{Multiple angles of a 3d scatterplot using the 3 most impactful
Variables}
\end{figure}

There is a clear distinction between the two family groups (Fig. 4)

\begin{longtable}[]{@{}
  >{\raggedright\arraybackslash}p{(\linewidth - 8\tabcolsep) * \real{0.4030}}
  >{\raggedleft\arraybackslash}p{(\linewidth - 8\tabcolsep) * \real{0.1343}}
  >{\raggedleft\arraybackslash}p{(\linewidth - 8\tabcolsep) * \real{0.2239}}
  >{\raggedleft\arraybackslash}p{(\linewidth - 8\tabcolsep) * \real{0.1194}}
  >{\raggedleft\arraybackslash}p{(\linewidth - 8\tabcolsep) * \real{0.1194}}@{}}
\caption{Logit of Family as a 5\% Ratio of Body Size, AIC =
122}\tabularnewline
\toprule\noalign{}
\begin{minipage}[b]{\linewidth}\raggedright
Variable
\end{minipage} & \begin{minipage}[b]{\linewidth}\raggedleft
Estimate
\end{minipage} & \begin{minipage}[b]{\linewidth}\raggedleft
Standard Error
\end{minipage} & \begin{minipage}[b]{\linewidth}\raggedleft
Z-value
\end{minipage} & \begin{minipage}[b]{\linewidth}\raggedleft
P-value
\end{minipage} \\
\midrule\noalign{}
\endfirsthead
\toprule\noalign{}
\begin{minipage}[b]{\linewidth}\raggedright
Variable
\end{minipage} & \begin{minipage}[b]{\linewidth}\raggedleft
Estimate
\end{minipage} & \begin{minipage}[b]{\linewidth}\raggedleft
Standard Error
\end{minipage} & \begin{minipage}[b]{\linewidth}\raggedleft
Z-value
\end{minipage} & \begin{minipage}[b]{\linewidth}\raggedleft
P-value
\end{minipage} \\
\midrule\noalign{}
\endhead
\bottomrule\noalign{}
\endlastfoot
(Intercept) & -3.715 & 3.21 & -1.16 & 0.248 \\
\texttt{Lower\ Leg\ Length\ Ratio} & 3.075 & 0.42 & 7.37 & 0.000 \\
\texttt{Lower\ Arm\ Length\ Ratio} & -6.562 & 0.83 & -7.94 & 0.000 \\
\texttt{Head\ Length\ Ratio} & 2.398 & 0.52 & 4.57 & 0.000 \\
\texttt{Eye\ Length\ Ratio} & -2.452 & 1.09 & -2.26 & 0.024 \\
\texttt{Naris-Snout\ Length\ Ratio} & -4.021 & 1.55 & -2.60 & 0.009 \\
\texttt{Snout-Vent\ Length} & -0.038 & 0.02 & -2.23 & 0.026 \\
\end{longtable}

Model Interpretations:

\begin{itemize}
\item
  Intercept: odds of being a tree frog when body ratios of Lower Leg
  Length, Lower Arm Length, Head Length, Eye Length, Naris-Snout Length,
  Snout-Vent Length are 0 are 0.024:1.
\item
  Lower Leg Length Ratio: When all other variables are held constant,
  for every 0.05 increase in lower leg length with respect to body size,
  the odds of being a tree frog change by 21.64 times.
\item
  Lower Arm Length Ratio: When all other variables are held constant,
  for every 0.05 increase in lower arm length with respect to body size,
  the odds of being a tree frog change by 0.001 times.
\item
  Head Length Ratio: When all other variables are held constant, for
  every 0.05 increase in head length with respect to body size, the odds
  of being a tree frog change by 11.00 times.
\item
  Eye Length Ratio: When all other variables are held constant, for
  every 0.05 increase in eye length with respect to body size, the odds
  of being a tree frog change by 0.086 times.
\item
  Naris-Snout Length Ratio: When all other variables are held constant,
  for every 0.05 increase in naris-snout length with respect to body
  size the odds of being a tree frog change by 0.017 times.
\item
  Snout-Vent Length: When all other variables are held constant, for
  every 1.0 increase in snout-vent length by 1 millimeter, the odds of
  being a tree frog change by 0.96 times
\end{itemize}

\textbf{Discussion}

Our goal of this project was to find the best morphological traits of
the order Anura to predict if a specimen would be a true toad or a tree
frog. From this statistical analysis we found that the best way to
predict a frog or a toad is by using 6 measurements: Body Size Ratios of
lower leg length, lower arm length, head length, eye length naris-snout
length, and an unratioed Snout-Vent Length. It is important to note that
the first 3 variables: lower leg length, lower arm length and head
length, are by far more important to the model than the last 3.
\emph{Each} of these three variables lowered the overal AIC (variability
explained by the model) by many hundreds, wheras the addition of the
latter 3 variables only lowered the AIC by a cumulative of 16.

Interestingly, though it would seem that having longer eye length or
naris snout length, with respect to body size, would raise the odds of
being a tree frog (Figs. 2 + 3), the model displays the opposite, as the
length increases the odds of being a tree frog is smaller. This is
likely due to the minimial variability in these variables, and the way
they interact with the other variables in the full model. If you have a
full specimen we reccomend trying to use all 6 variables, but in a
pinch, the measurements of lower leg, lower arm, and head length
relative to body size will suffice.

Though we used a step-wise exploratory method for arriving on our model,
rather than testing specific variables of interest, the fact that
measurements of lower leg length and lower arm length make a lot of
sense in the context of the data. As discussed earlier, one of the most
unique features to the tree frogs are their strong jumping legs. The
fact that the odds of being a tree frog change by 21.64 times for every
.05 increase in relative lower leg length strongly affirms the claim by
Paris (2000). Our model also stated that the odds of being a tree frog
change by .001 times for every .05 increase in lower arm length as a
ratio to body size. This means that the trait of longer lower arms is
highly specialized to toads, which strongly corroborates the connection
made by Sanchez et. al (2010) of soil hardness and burrowing habits, and
is a clear distinction between the two families.

One potential confounding variable to our data is that we do not know
how these specimens were stored. The authors did not clairfy if
specimins were stored in jars or as skeletons, this difference could
warp measurements due to layers of skin and muscle, especially in
smaller variables, such as naris-snout length. Furthermore, preserved
specemins are known to shrink over time, despite the best preservation
techniques, specimins of different ages could have shrunken or become
deformed at different rates, warping them from their `true' measurement,
adding extra variability to the dataset.

This model can be generalized to identify frogs that fall in the family
of Hylidae or Bufonidae, which limits its ability to identify an unknown
frog specemin, though if it is known to fall into one of those two
families, it could be especially useful then. More complex models could
also be made to narrow down a frog from a wider selction of families or
species. The source of the data library of specemins is proudly boasted
to be `world wide' though our spececific data is not connected to a
location, so it is hard to say where in the world we can generalize this
data to. If we assume the sampling of species and locations is random
though, it could be feasible to generalize it to any species within
these two families around the world.

One of the motivations behind our research was the declining population
of frogs, which is still vastly misunderstood. With this analysis, we
have identified some clear defining characteristcs to make the
identification of diceased frogs easier. Furthermore, our model found
clear evidence of family wide nich partioning in Tree frogs and True
Toads. Tree frogs lower \emph{legs} are longer and specialized for
jumping between trees, while true toads lower \emph{arms} are longer and
specialized for digging burrows on the forest floor. Thus this model
could be used to aid the conservation of frogs in a 2 fold manor. First,
a specimen could be identified by a series of body measurements, and
then linked to a location and larger dataset. Secondly, by identifying
which traits are important in these specimins, we can begin to look for
problems that relate to their specialization, such as in deforestatin or
soil quality. This analysis could be done via a hot spot analysis, where
a subject of interest is identified and then algorithmically analyzed to
reveal spatial trends correlated and shows evidence of a phenomenon
happening with a confidence level.

Future logistic models should work to assess differentiations of other
families, and potentially explore specimens on a species level.

\subsection{Sources:}\label{sources}

AmphibiaWeb---Hylidae. (n.d.). Retrieved May 16, 2025, from
\url{https://amphibiaweb.org/lists/Hylidae.shtml}

Andreas, M., Reiter, A., Cepáková, E., \& Uhrin, M. (2013). Body size as
an important factor determining trophic niche partitioning in three
syntopic rhinolophid bat species. Biologia, 68(1), 170--175.
\url{https://doi.org/10.2478/s11756-012-0139-1}

Bufonidae (bufonidés, Bufonids, crapauds, Toads). (n.d.). Animal
Diversity Web. Retrieved May 16, 2025, from
\url{https://animaldiversity.org/accounts/Bufonidae/}

Colloca, F., Carpentieri, P., Balestri, E., \& Ardizzone, G. (2010).
Food resource partitioning in a Mediterranean demersal fish assemblage:
The effect of body size and niche width. Marine Biology, 157(3),
565--574. \url{https://doi.org/10.1007/s00227-009-1342-7}

Parris, K. M. (2001). Distribution, habitat requirements and
conservation of the cascade treefrog (Litoria pearsoniana, Anura:
Hylidae). Biological Conservation, 99(3), 285--292.
\url{https://doi.org/10.1016/S0006-3207(00)00194-4}

Sanchez, L. C., Busch, M., \& Madanes, N. (2010). Terrestrial habitat
use by the burrowing toad, Rhinella fernandezae (Anura: Bufonidae).
Folia Zoologica, 59(2), 122--128.
\url{https://doi.org/10.25225/fozo.v59.i2.a7.2010}

Wake, D. B. (1991). Declining Amphibian Populations. Science, 253(5022),
860--860. \url{https://doi.org/10.1126/science.253.5022.860}

\end{document}
